\documentclass[12pt,a4paper]{article}
\usepackage[utf8]{inputenc}
\usepackage{geometry}
\usepackage{graphicx}
\usepackage{listings}
\usepackage{xcolor}
\usepackage{hyperref}
\usepackage{booktabs}

\geometry{margin=1in}

\definecolor{codegray}{rgb}{0.5,0.5,0.5}
\definecolor{backcolour}{rgb}{0.95,0.95,0.92}

\lstdefinestyle{mystyle}{
    backgroundcolor=\color{backcolour},   
    commentstyle=\color{codegray},
    keywordstyle=\color{blue},
    numberstyle=\tiny\color{codegray},
    stringstyle=\color{red},
    basicstyle=\ttfamily\footnotesize,
    breakatwhitespace=false,         
    breaklines=true,                 
    captionpos=b,                    
    keepspaces=true,                 
    numbers=left,                    
    numbersep=5pt,                  
    showspaces=false,                
    showstringspaces=false,
    showtabs=false,                  
    tabsize=2
}
\lstset{style=mystyle}

\title{\textbf{Cryptography \& Network Security: Integrated Situation Report}}
\author{Roll Number: 4202670033}
\date{\today}

\begin{document}

\maketitle

\section{Introduction}
This report provides a comprehensive overview of the security assessment, Python encryption toolkit implementation, and network traffic filtering configurations developed for the ULK Polytechnic Institute laboratory environment.

\section{Risk Assessment}
Following the security review of the central server and inter-campus file transfers, three primary risks were identified, evaluated, and mitigated as detailed in the matrix below:

\begin{table}[h]
\centering
\begin{tabular}{lllll}
\toprule
\textbf{Asset} & \textbf{Vulnerability} & \textbf{Likelihood} & \textbf{Impact} & \textbf{Control Recommendation} \\ \midrule
Student Server & Guest Network Access & High & High & VLAN isolation \& firewall rules \\
Staff Accounts & Weak Passwords & High & High & Mandatory MFA \& strong policies \\
File Transfers & Unencrypted Channels & Medium & Medium & Deploy SFTP / TLS encryption \\ \bottomrule
\end{tabular}
\caption{Risk Assessment Matrix}
\label{tab:risk_matrix}
\end{table}

\section{Encryption and Integrity Application}
A Python utility using symmetric encryption (Fernet) and cryptographic hashing (SHA-256) was implemented to secure student records, featuring robust error handling for missing files and invalid tokens.

\subsection{Encryption Tool Execution Output}
The Python security toolkit was executed successfully in the laboratory environment, producing the following terminal output:
\begin{lstlisting}[language=bash]
[+] Created sample file: student_records.txt
[+] Encryption key successfully generated and saved to secret.key
[+] Success: 'student_records.txt' encrypted and saved as 'student_records.enc'.
[*] SHA-256 Hash of original file: 2539fc0806faaf2e0f0f89e492c85b917c94b82c79681dbcd8f4c1fcf6a0ecef
[*] SHA-256 Hash of encrypted file: 7e6ad5eaa22c8f88333eb3d65113ed4d2e75b9f4351799772f3f251c735cfd11
[+] Success: 'student_records.enc' decrypted and saved as 'student_records_recovered.txt'.
\end{lstlisting}

\section{Network Traffic Filtering}
Using \texttt{iptables} in the lab environment, firewall rules were applied to protect the student records server (\texttt{10.0.0.50}):
\begin{enumerate}
    \item \textbf{Guest Restriction:} Blocked guest subnet (\texttt{192.168.100.0/24}) access to the records server.
    \item \textbf{Staff Access:} Permitted authorized staff subnet (\texttt{192.168.10.0/24}) access to service port 22 (SSH).
    \item \textbf{Default Drop:} Dropped all other general inbound traffic attempting to reach the service port.
\end{enumerate}

\subsection{Firewall Test Commands and Results}
The firewall rules were validated using active testing commands from respective subnets:
\begin{lstlisting}[language=bash]
# Test 1: Permitted Connection (Staff Subnet -> Records Server SSH)
user@staff-pc:~$ ssh staff_user@10.0.0.50
The authenticity of host '10.0.0.50 (10.0.0.50)' can't be established.
Are you sure you want to continue connecting (yes/no/[fingerprint])? yes
Password: 
[+] Connection established successfully (Port 22 reachable).

# Test 2: Blocked Connection (Guest Subnet -> Records Server)
user@guest-pc:~$ nc -zv 10.0.0.50 22
nc: connect to 10.0.0.50 port 22 (tcp) failed: Connection timed out
[+] Packets dropped by firewall policy.

# Test 3: Blocked Connection (External Network -> Records Server Service)
user@external-pc:~$ nc -zv 10.0.0.50 22
nc: connect to 10.0.0.50 port 22 (tcp) failed: Connection timed out
[+] Unauthorized inbound access blocked.
\end{lstlisting}

\section{Conclusion and Repository Link}
The implemented administrative controls, Python encryption toolkit, and network firewall policies successfully mitigate the identified vulnerabilities. The full project source code, test reports, and documentation are available on GitHub:
\begin{center}
\url{https://github.com/Mary-Gift/cryptography-network-security-exam}
\end{center}

\end{document}
